%----------------------------------------------------------------------------------------
%	PACKAGES AND THEMES
%----------------------------------------------------------------------------------------
\documentclass{beamer}
% Load a bunch of useful packages:
\usepackage{amssymb,amsmath,amsfonts,mathtools} % useful math fonts and symbols
\usepackage{geometry} % allows changing margins and sizes of stuff
\usepackage{hyperref} % allows referencing of lines of text, urls, figures etc
\usepackage{natbib} % allows citation referencing from a .bib file
\bibliographystyle{apalike} % American Psychological Association style guide
\usepackage{graphicx} % to include graphics from the figures folder with helpful formating options
\usepackage{tikz}
\usetikzlibrary{arrows}
\usepackage{pgfplots}
\usetikzlibrary{calc}
%---------------------------------------------%
\usepackage{color} % custom color definitions:%
    %UOregon colors:                          %
    \definecolor{UOGreen}{RGB}{18, 71, 52}    %
    \definecolor{UOYellow}{RGB}{254, 225, 35} %
    %Secondary official colors:               %
    \definecolor{LegacyGreen}{HTML}{104735}   %
    \definecolor{GrassGreen}{HTML}{489D46}    %
    \definecolor{LimeGreen}{HTML}{8ABB40}     %
    \definecolor{Chartreuse}{HTML}{E2E11B}    %
    \definecolor{Berry}{HTML}{8D1D58}         %
    \definecolor{DarkBlue}{HTML}{004F6E}      %
    \definecolor{LightBlue}{HTML}{00A5B5}     %
    \definecolor{crimson}{RGB}{ 170, 4, 36 }
    \definecolor{darkblue}{RGB}{ 4, 47, 170 }
    \definecolor{brown}{RGB}{ 111, 71, 2 }
    \definecolor{periwinkle}{RGB}{ 90, 177, 204 }
    \definecolor{ducksgreen}{HTML}{007030}
%----------------------------------------------

\usepackage{setspace} % to specify single, double or one-half spacing of text
\usepackage{indentfirst} % indent first line of a text paragraph
\usepackage{multicol} %multipage ability
\usepackage{multirow}
\usepackage{ulem} % \ul (underline) command which will break over line ends
\usepackage{amsthm} % allows standarized theorem commands
\setbeamertemplate{theorems}[numbered]
\theoremstyle{plain}
\newtheorem{assume}{Assumption}
\newtheorem{define}{Definition}
\usepackage{breqn} %automatic line breaking
\normalem

% Theming and Appearance Setup:
\usetheme{metropolis}
\metroset{block=fill}
%\usetheme{default}
\usefonttheme{professionalfonts}
\fontfamily{ppl}\selectfont
\usecolortheme[named=UOGreen]{structure}
\setbeamertemplate{footline}[frame number]
\setbeamertemplate{headline}{} %Removes Section Index on top of slide
\beamertemplatenavigationsymbolsempty
\hypersetup{
    colorlinks=true,
    linkcolor=DarkBlue,
    filecolor=Berry,
    citecolor=GrassGreen,
    urlcolor={LightBlue},
    pdftitle={Trade and Labor Market Dynamics},
    pdfpagemode=FullScreen
}

%------------------------------------------------------------------------------%TITLE PAGE
%------------------------------------------------------------------------------
% The title
\title{}
\author{Dante Yasui }
\institute{EC327 Game Theory}
\date{Winter 2024}
\titlegraphic{\includegraphics[scale=.4]{UOSignature-356.png}}


%-----------------------------------------------------------------------------
%	PRESENTATION SLIDES
%-----------------------------------------------------------------------------

\begin{document}

\begin{frame}[plain]
    % Print the title page as the first slide
    \titlepage
\end{frame}
\addtocounter{framenumber}{-1}

\begin{frame}[plain]{Outline}
  \tableofcontents
\end{frame}
\addtocounter{framenumber}{-1}

% - - - - - - - - - - - - - - - - - - - - - - - - - - - - - - - - - - - - - -

\begin{frame}{}
  \begin{itemize}
    \item We have already talked about how different types of games 
    favor some agents over others 
    \begin{itemize}
      \item I.e., first mover advantage, 
      \item second mover advantage, 
      \item asymmetric info 
    \end{itemize}
    \item So it would make sense that if players can \textit{manipulate}
    the rules of a game in their favor, they will try to do so.
  \end{itemize}  
\end{frame}

% - - - - - - - - - - - - - - - - - - - - - - - - - - - - - - - - - - - - - -

\begin{frame}{}
  \begin{itemize}
    \item We can think about adding a first-stage to our original game 
    \begin{itemize}
      \item \textbf{First Stage:} specify how you will act in second stage 
      \item \textbf{Second Stage:} the original game
      \begin{itemize}
        \item but now players set their beliefs based on what happened
        in the first stage.
      \end{itemize}
    \end{itemize}
  \end{itemize}  
\end{frame}

% - - - - - - - - - - - - - - - - - - - - - - - - - - - - - - - - - - - - - -

\begin{frame}{}
  \begin{itemize}
    \item Different first stage actions correspond to what we will call:
    \begin{itemize}
      \item \alert{commitments},
      \item \alert{threats},
      \item or \alert{promises}
    \end{itemize}
    \item Whether any of these actions is \textbf{effective}
    depends on the beliefs of the other player(s).
    \begin{itemize}
      \item The \alert{credibility} of a strategic move \textit{matters}.
    \end{itemize}
  \end{itemize} 
\end{frame}

% - - - - - - - - - - - - - - - - - - - - - - - - - - - - - - - - - - - - - -

\begin{frame}{}
  Examples of Strategic Moves:
  \begin{itemize}
    \item Amazon publicly \textit{commits} to going carbon neutral by 2040
    \item Parents \textit{promise} "you will get a PS5 if you get all A's" 
    to their children
    \item Nuclear powers \textit{threaten} "Mutually Assured Destruction"
    to each other in \alert{brinksmanship} games
  \end{itemize}
\end{frame}

% - - - - - - - - - - - - - - - - - - - - - - - - - - - - - - - - - - - - - -

\begin{frame}{}
  Conditional Strategic Moves 
  \begin{itemize}
    \item I might declare a \textbf{response rule} 
    which is a move that depends on someone else's behavior
    \item I might try to take action to stop someone from doing something 
    through a \textbf{deterrence} strategy
    \item Or I could try to get someone to do something 
    through \textbf{compellence}
    \item \texbf{deterrence} or \textbf{compellence} could take the form 
    of a \alert{threat} or \alert{promise}:
    \begin{itemize}
      \item A \alert{threat}: 
      "Unless your action conforms to what I want, then I will \textit{harm} you"
      \item A \alert{promise}:
      "If your action conforms to what I want, I will \alert{reward} you"
    \end{itemize}
  \end{itemize}
\end{frame}

\section{Classification of Strategic Moves}

\section{Commitments}

\section{Threats and Promises}

\section{Additional Topics}

\end{document}
